\section{Q1：如何让 AI 判断什么研究值得做？}
\label{sec:taste}

Research taste 关心的是结果尚未确定时的决策：哪些未知值得研究，哪些提案值得投入资源，以及什么结果应该改变下一步实验。因此，评价必须说明决策情境。对已完成论文排序、筛选未执行提案和选择下一个实验彼此相关，但可见信息与反馈不同。本文将新颖性、可行性、预期信息增益和资源需求分开讨论，不预设某个单一分数就能表示 scientific value。

\subsection{不同反馈究竟测量了什么}

AI Can Learn Scientific Taste 将学术共同体的反响转化为训练信号：在同领域、同年份论文中依据引用差异构造偏好，训练 Scientific Judge，再用它奖励 Scientific Thinker。作者、机构和主题控制排除了一部分明显的捷径。不过，生成器仍通过 LLM 对潜在影响力的比较进行评价，并未执行提案 \cite{ai-can-learn-scientific-taste}。这支持偏好能够被学习并传递到生成端，但尚未确认这种偏好与实际科研价值的关系。

HindSight 的目标不同：将生成想法与后续论文匹配，再利用发表相关的影响信号加权。它使回顾式比较成为可能，却也奖励与后来已发表研究的一致性；没有后续匹配的方向未必没有价值。更关键的是，文中模型知识截止时间晚于检索截止时间，仅限制检索文献不能排除模型知道评价窗口中的内容 \cite{hindsight}。因此，本文将其视为 retrospective alignment 的证据，而非严格的 prospective discovery forecast。

TasteVal 从论文偏好转向实验决策：给定研究问题，固定 coding agent，衡量实验算力与研究进展的关系。这有助于区分规划、解释与代码实现，但测量仍依赖 coder、任务、预算和评分规则，而且任务未公开发布 \cite{tasteval}。其结果不能直接证明开放式选题能力，也不能将这一协议当作已可获取的公共训练集。

执行还可能改变评价本身。Ideation--Execution Gap 在研究者选择的主题内随机分配人类或 LLM 提案，再对执行成果进行专家盲审；LLM 提案在最初评价中的优势经实施后减弱 \cite{ideation-execution-gap}。它比单纯提案评分更接近实际成果，但人类执行和特定 NLP 情境仍是实验条件的一部分。这些研究未必相互矛盾：预判共同体反响、获得评审认可、选择有效实验，本来就是不同的预测目标。

\begin{table}[t]
\centering
\small
\caption{Research taste 的反馈来源及其支持的推断；任何一种目标都不能单独定义科学价值。}
\label{tab:taste-feedback}
\begin{tabularx}{\linewidth}{@{}P{0.22\linewidth}YY@{}}
\toprule
反馈来源 & 可观察目标 & 尚未解决的问题 \\
\midrule
共同体反响 \cite{ai-can-learn-scientific-taste,hindsight} & 引用偏好或后续发表匹配 & 未发表的替代方向及因果意义上的科学价值 \\
专家或证据判断 \cite{rinobench,opennovelty,innoeval} & 新颖性判断与评审一致性 & 证据缺失与评审偏好 \\
受控扰动 \cite{scistylebench} & 对风格稳定、对内容变化敏感 & 稳健性是否预示有效研究 \\
实际执行 \cite{tasteval,ideation-execution-gap} & 预算内进展或执行后评价 & 对实现方式与任务选择的依赖 \\
\bottomrule
\end{tabularx}
\end{table}

\subsection{有证据的判断与对捷径的抵抗}

一条路线是缩小判断范围并公开依据。RINoBench 对照从专家评审中提取的标签评价新颖性分数和理由；OpenNovelty 将论文贡献主张与相关工作逐项比较；InnoEval 则结合检索证据与多种评审视角 \cite{rinobench,opennovelty,innoeval}。这些方法让判断更容易检查，但仍依赖检索覆盖和比较的可靠性。看似合理的理由可以伴随错误分数；证明主张区别于已有工作，也不能证明它值得研究。

SciStyleBench 提出互补问题：judge 对科学内容敏感，还是对其表达方式敏感？它将风格变化与内容削弱、增强分别干预。关键是同时满足两项要求：不受无关修辞影响，又能识别方法变化；恒定评分器虽满足不变性，却没有判断价值 \cite{scistylebench}。本文据此认为，evidence grounding 与 shortcut resistance 应联合评价。增加文献可以补充事实，但不会自动让决策规则忠实于这些事实。

\subsection{局限与开放问题}

现有证据支持多种 taste 的操作化定义，尚未建立统一能力。廉价标签利于扩展，昂贵执行提供更直接但范围更窄的反馈；两者都不能消除评价分歧或替研究者决定目标。如何奖励有信息量的负结果、何时允许 judge 弃权、收益能否跨任务和独立实现保留，仍是开放问题。前瞻研究应在观察结果前冻结提案、评审器、预算和可访问信息。以下方案均为拟议实验，并非已证实的结论。

\subsection{拟议计划一：结合执行结果与主动标注的提案排序}

\paragraph{研究假设与方法。}
在固定标注预算下，结合已完成 pilot 的结果与主动请求的专家偏好，比仅学习执行前评分更能选出可复现、具有信息量的研究。以小型 pairwise judge 为起点，优先请求不确定或存在分歧的样本对，同时保留随机标注比例并记录抽样概率。专家依据预先定义的 rubric 比较训练阶段已执行 pilot，承认有效负结果的价值，并区分实现失败与假设被否定。

\paragraph{任务与划分。}
围绕冻结小模型的校准、检索评分和鲁棒性构造十二个低成本 AI 研究任务，每个包含六项候选干预。六个任务用于训练、三个用于开发、三个用于封存评价；共享数据集和代码模板的任务必须放在同一分区。评价阶段以相同 coder 和预算执行全部候选，使用独立随机种子复现；排序冻结前不开放结果。这个适度规模的设计检验跨任务迁移，不声称测量普遍科研品味。

\paragraph{对照、指标与风险。}
比较随机选择、直接提示 judge、仅提案偏好训练，以及随机标注的执行结果训练；分别消融主动选择与执行反馈。指标包括 top-$k$ 中可复现且有信息量结果的比例、相对全部候选执行结果的 ranking regret、校准、标注成本和实验成本，并报告跨任务不确定性。主要风险是结果噪声、偏爱容易项目，以及主动抽样遗漏陌生想法；随机样本提供审计入口，但不能消除这些风险。

\subsection{拟议计划二：抵抗捷径但保留内容辨别力的 judge}

\paragraph{研究假设与方法。}
将主张、支持证据与表达线索分开，有望提高风格变化下的排序稳定性，同时保留对无效方法的敏感性。围绕约一百个提案构造受控变体：一组改变措辞、确定语气和声望线索；另一组移除必要对照、引入矛盾或缺乏依据的机制。由独立标注者核实干预是否保留或改变目标内容，再用不变性与内容辨别两个目标训练轻量 pairwise evaluator。

\paragraph{任务、对照与评价。}
按原始提案、主题和改写模型划分，确保同一想法的所有变体属于同一分区。比较直接提示、检索增强评价、文本规范化，以及两个训练目标分别单独使用的结果。测量风格导致的分数变化、实质缺陷识别率和对抗表达下的 top-$k$ 更替，并在计划一封存的已执行任务上检验稳健性是否改善选择质量。

\paragraph{失败风险。}
改写可能改变语义，规范化可能删掉有用细节，人工缺陷也可能比真实科研弱点容易识别，因此必须人工审计并保留未经改写的提案。即使通常的 judge agreement 上升，只要风格敏感性下降伴随缺陷辨别或下游选择质量退化，就应判为失败。
